\section{Glosario}
\noindent
\begin{description}
Glosario \textbf{ordenado} de términos cuyo significado se ha buscado en algun momento, o se desconocía \textit{a priori} o en general poco frecuente y obscuro
    
% A-M
    \item[Constitución] Norma fundamental del estado. 
    \item[Colegio] Sociedad o corporación de personas de una misma profesión, a la que generalmente se atribuyen funciones de ordenación y disciplina de la actividad profesional.
    \item[Derogar] Dejar sin efecto una norma vigente.
    \item[Disposición] Precepto legal o reglamentario, deliberación, orden y mandato de la autoridad.
    \item[Impugnar] Combatir, contradecir, refutar
    \item[Impugnar] Interponer un recurso contra una resolución judicial 
    \item[Interdecir] Vedar, prohibir
    \item[Interponer] Dar inicio a un recurso legal
% N-Z
    \item[Ley] Precepto dictado por la autoridad competente, en que se manda o prohíbe algo en consonancia con la justicia y para el bien de los gobernados.
    \item[Obscuro] Confuso, poco inteligible
    \item[Ordenamiento Juridico] Conjunto de todas las leyes y normas que rigen en un lugar y momento específico.
    \item[Retroactividad] Extensión de la aplicación de una norma a hechos y situaciones anteriores a su entrada en vigor o a actos y negocios jurídicos.
    \item[Texto articulado] Texto que tiene artículos con la finalidad de crear una nueva norma sobre una determinada materia
    \item[Texto refundido] Texto que recoge y unifica varias leyes que regulan la misma materia
\end{description}
